\documentclass[11pt,a4paper]{article}
\usepackage[utf8]{inputenc}
\usepackage{graphicx}
\usepackage[left=2.5cm,top=3cm,right=2.5cm,bottom=3cm,bindingoffset=0.5cm]{geometry}
\usepackage{macros/AEDLogica, macros/AEDEspecificacion, macros/AEDTADs}
\usepackage{caratula}


\titulo{Trabajo práctico}
\subtitulo{AirMiles}

\fecha{\today}

\materia{Algoritmos y Estructuras de Datos}
\grupo{Nombre del grupo} 

\integrante{Bravo, Leonardo}{217/24}{bravopleo@gmail.com}
\integrante{Garcia, Lautaro}{1104/24}{lautarogonzalogarcia@outlook.com}
\integrante{Tapia, Avril}{1420/25}{celeste.avril.023@gmail.com}
\integrante{Gómez, Héctor}{921/23}{hectorluisgomezmoya@gmail.com}


% Declaramos donde van a estar las figuras
% No es obligatorio, pero suele ser comodo
\graphicspath{{../static/}}

% Asi pueden escribir nuevos comandos. 
% Este por ejemplo asegura q los nombres 
% que figuren con una tipografia diferenciada  
\newcommand{\Tipo}[1]{\mathsf{#1}} 
% la sintaxis es \newcommand{\nombreDeLaMacro}[cantidadDeParametros]
%{Lo que va ser remplazado por el macro} 
\newcommand{\norm}[1]{\vert #1\vert}



\begin{document}

% https://docs.google.com/spreadsheets/d/1p3TdUczBpR4ySYpvJCpe2JbsCuC5z6lokaHnCKXhCpI/edit?gid=0#gid=0


\maketitle
\section{Renombres de tipos y notas}

\subsection{Tipos}

    \ES{Id}{\Z}
    \ES{Millas}{\R} 
    \ES{Factor}{\R}
    \ES{Tope}{\R}
    \ES{Codigo}{\str}
    \ES{TieneTope}{\bool}
    %se, parece que no esta definido, pero con que mandemos un text o algo antes que diga " es enum" y esa secuencia de valores, ya fue
    % ahi ta . voy a seguir agregando las funciones que no marcaron
    \ES{Categoria}{enum \langle BRONCE, PLATA, ORO, PLATINO\rangle}
    % \renombre{TipoTransaccion}{\Tipo{enum}\ \langle EMISOR, RECEPTOR, VUELO, CANJE, REFERIDO \rangle} \\
    \ES{TipoTransaccion}{enum \langle EMISOR, RECEPTOR, VUELO, CANJE, REFERIDO\rangle}
    % nos faltan datos, tenemos que representar la distancia y el multiplicador del codigo
    % att, podemos representar solo el mult de codigo y usar millas como la "distancia"
    % y con un aux representamos los puntos de millas acumuladas, es decir, otro aux que te de el calculo de la multiplicacion o en el mismo aux que tenemos
    \ES{Transaccion}{\struct{tipo: \Tipo{TipoTransaccion},\hspace{.1cm} distancia: \Tipo{Millas},\hspace{.1cm} millas: \Tipo{Millas}, \\ \tab\tab\tab\tab factorPromo: \Tipo{Factor},\hspace{.1cm} codigo: \Tipo{Codigo},\hspace{.1cm}\ contraparte: \Tipo{Id}}} 
    \ES{Historial}{\seq{\Tipo{Transaccion} }}

\subsection{Notas}

    \begin{itemize}
        
        \item $\oplus$ ($\lor$ exclusivo): sean $a$ y $b$ proposiciones lógicas, $a \oplus b ~ = ~ (a \land \neg b) \lor (\neg a \land b)$

        
        %\item En el desarrollo de la implementación del TAD a continuación, en base a la información brindada del mismo en %el enunciado del trabajo práctico, consideramos no necesario guardar el registro de cierta información, ya que si %el TAD requiriese de la misma para desarrollar sus procedimientos y/u operaciones la tendriamos en cuenta y la %registraríamos en alguno de nuestros observadores.
        

        \item \textbf{Contraparte de una transacción.} Cada transacción del historial de un usuario guarda en el
          campo $contraparte$ el $\Tipo{Id}$ del otro usuario involucrado en ella:
        \begin{itemize}
              \item si $t.tipo = \Tipo{EMISOR}$, $t.contraparte$ es el usuario que \emph{recibió} las millas;
              \item si $t.tipo = \Tipo{RECEPTOR}$, $t.contraparte$ es el usuario que \emph{envió} las millas;
              \item si $t.tipo = \Tipo{REFERIDO}$, $t.contraparte$ es el otro usuario de la referencia
                    (el referidor para el usuario nuevo, y el usuario nuevo para el referidor);
              \item si $t.tipo = \Tipo{VUELO}$ o $t.tipo = \Tipo{CANJE}$ no hay otro usuario involucrado, y se usa
                    $t.contraparte = sinContraparte$, donde $sinContraparte = 0$ (ver auxiliar
                    \text{sinContraparte}).
          \end{itemize}
          El valor $0$ nunca es el $\Tipo{Id}$ de un usuario: \text{registrarUsuario} y
          \text{registrarUsuarioConReferidor} solo crean usuarios con $\Tipo{Id} > 0$, y son los únicos
          procedimientos que agregan claves a $usuarios$. Por lo tanto, en todo estado del TAD vale
            $sinContraparte \notin claves(usuarios)$, y $t.contraparte = sinContraparte$ significa
          inequívocamente que la transacción $t$ no tiene contraparte.
    \end{itemize}



\section{AirMiles}


\begin{tad}{AirMiles}

\subsection{Observadores}
% CORRECCIÓN: "No es la forma de escribir un struct"
\obs{usuarios}{\dict{\Tipo{Id}}{\Tipo{Historial}}} \\
\obs{promos}{\dict{\Tipo{Codigo}}{\Tipo{struct}\langle \Tipo{factor}:Factor,\hspace{.1cm} tope: \Tipo{Tope},\hspace{.1cm} tieneTope: \Tipo{TieneTope}\rangle}}


\subsection{Auxiliares}

% Valor reservado para el campo contraparte de las transacciones sin otro usuario involucrado
% (VUELO y CANJE). Ningún usuario tiene este Id: ver Notas.
\aux{sinContraparte}{}{\Tipo{Id}}{0}

\bigskip

% CORRECCIÓN: "Los enum no son strings, no van entre comillas. emisor != emision."
\auxLargo{millasDisponibles}{h: \Tipo{Historial}}{\R}{
    \sum_{k=0}^{|h|-1} \Tipo{IfThenElse}(h[k].tipo = \comillas{EMISOR} \hspace{.1cm} \lor \hspace{.1cm} h[k].tipo = \comillas{CANJE},
    - h[k].millas,
    h[k].millas)
}

\bigskip

\auxLargo{millasAcumuladas}{h: \Tipo{Historial}}{\R}{
    \sum_{k=0}^{|h|-1} \IfThenElse{h[k].transaccion = EMISOR \hspace{.1cm} \lor \hspace{.1cm} h[k].transaccion = CANJE}{0}{h[k].millas}
}

\bigskip

\auxLargo{millasCanjeadas}{h: \Tipo{Historial}}{\R}{
    \sum_{k=0}^{|h|-1} \IfThenElse{h[k].transaccion = CANJE}{h[k].millas}{0}
}

\bigskip

\auxLargo{totalMillasPorVuelo}{h: \Tipo{Historial}}{\R}{
    \sum_{i=0}^{|h|-1} \IfThenElse{h[i].transaccion = VUELO}{h[i].millas}{0}
}

\bigskip

\auxLargo{totalMillasCanjeadas}{h: \Tipo{Historial}}{\R}{ 
    \sum_{i=0}^{|h|-1} \IfThenElse{h[i].transaccion = CANJE}{h[i].millas}{0}
}

\bigskip

% Antes era totalMillasEmitidasRecibidas y sumaba emitidas y recibidas juntas (se perdía información).
% Ahora se separa en dos auxiliares: uno para las emitidas y otro para las recibidas.
\auxLargo{totalMillasEmitidas}{h: \Tipo{Historial}}{\R}{
    \sum_{i=0}^{|h|-1} \IfThenElse{h[i].transaccion = \Tipo{EMISOR}}{h[i].millas}{0}
}

\bigskip

\auxLargo{totalMillasRecibidas}{h: \Tipo{Historial}}{\R}{
    \sum_{i=0}^{|h|-1} \IfThenElse{h[i].transaccion = \Tipo{RECEPTOR}}{h[i].millas}{0}
}
\bigskip

\auxLargo{millasGanadasPorVuelo}{distancia: \R, historial: \Tipo{Historial}, factor: \Tipo{Factor}}{\R}{
    distancia \cdot factorCategoriaUsuario(historial) \cdot factor
}

\bigskip

\auxLargo{factorCategoriaUsuario}{h: \Tipo{Historial}}{\R}{
    (1 + factorCondicional(millasAcumuladas(h), 0.5, 0, 10000) + \\ \tab factorCondicional(millasAcumuladas(h), 0.5, 0, 50000))
    \cdot \\ \tab\tab factorCondicional(millasAcumuladas(h), 0, 1, 1000000)
}

\bigskip

\auxLargo{factorCondicional}{m: \Tipo{Millas}, factor1: \R, factor2: \R, rango: \mathbb{N}}{\R}{
    \IfThenElse{m \geq rango}{factor1}{factor2}
}
\bigskip

\auxLargo{sumaDeMillas}{h: \Tipo{Historial}, c: \conj{Transaccion}}{\R}{
    \sum_{i=0}^{|h|-1} \IfThenElse{h[i].transaccion \in c}{h[i].monto}{0}
}

\bigskip

\subsection{Predicados}

    \predLargo{esRankingOrdenado}{
        ranking:\seq{\Tipo{Id}}, 
        usuarios: \dict{\Tipo{Id}}{\Tipo{Historial}}
        }
        {
            \paraTodo{\,i}{\Z}{0 \leq i < |ranking|-1
            \implicaLuego \\ (ranking[i] \in claves(usuarios) ~ \land ~ ranking[i+1] \in claves(usuarios)) \yLuego \\ millasDisponibles(usuarios[ranking[i]]) \geq millasDisponibles(usuarios[ranking[i+1]])}
        } 

\bigskip

    \predLargo{tieneALosTopN}{
        ranking:\seq{\Tipo{Id}},
        usuarios: \dict{\Tipo{Id}}{\Tipo{Historial}}
        }
        {
            \paraTodo{\,i,j}{\Tipo{Id}}{(i \in claves(usuarios) \ \land \ i \notin ranking) \ \land \ (j \in claves(usuarios) \ \land \ j \in ranking) \implicaLuego millasDisponibles(usuarios[ranking[i]]) < \\ millasDisponibles(usuarios[ranking[j]])}
        } 

\bigskip

    \predLargo{esPromoValida}{\\
        \tab promo: \Tipo{struct}\langle factor: \Tipo{Factor}, tope: \Tipo{Tope}, tieneTope: \Tipo{TieneTope} \rangle, \\
        \tab millasAGanar: \Tipo{Millas} \\}{
        (\lnot promo.tieneTope \implica promo.Tope = 0) \land \\
            \tab (promo.tieneTope \implica millasAGanar \leq promo.Tope) \\
    }

\bigskip

    \predLargo{sinRepetidos}{   
        ranking: \seq{\Tipo{Id}}
        }
        {
            \paraTodo{i,j}{\Z}{0 \leq i < |ranking| ~ \land ~ 0 \leq j < |ranking| ~ \implicaLuego ((ranking[i] = ranking[j]) \iff i = j)}
        } 

\bigskip

    \predLargo{sonTodosUsuarios}{
        ranking: \seq{\Tipo{Id}},
        usuarios: \dict{\Tipo{Id}}{\struct{\Tipo{Historial},\Tipo{EMISOR}}}       
        }   
        {
            \paraTodo{i}{\Z}{0 \le i < |ranking| \implicaLuego (ranking[i] \in claves(usuarios)) }
        } 

\bigskip   


    \predLargo{sinParesRepetidos}{
        tuplas: \seq{\tupla{\Tipo{Id},\Tipo{Id}}}
        }
    {    \paraTodo{i,j}{\Z}{0 \leq i < |tuplas| \land 0 \leq j < |tuplas| \implicaLuego \\
                \tab\tab ((tuplas[i]_0 = tuplas[j]_0 \land tuplas[i]_1 = tuplas[j]_1) \iff i = j)
        }
    }

\bigskip

    % Ahora se calculan a partir del historial (campo contraparte) en lugar de la secuencia emision.
    \pred{sonExclusivos}{h_1: \Tipo{historial}, h_2: \Tipo{historial}, u_1: \Tipo{id}, u_2: \Tipo{id}}{\\
        \tab\tab soloEmitioA(h_1, u_2) \land soloEmitioA(h_2, u_1)
    }

\bigskip

    % El usuario dueño de h emitió al menos una transferencia, y todas las que emitió fueron hacia u.
    \predLargo{soloEmitioA}{h: \Tipo{historial}, u: \Tipo{id}}{ 
        \existe{i}{\Z}{0 \leq i < |h| \yLuego h[i].tipo = \Tipo{emisor}} \land \\
        \paraTodo{i}{\Z}{(0 \leq i < |h| \yLuego h[i].tipo = \Tipo{emisor}) \implicaLuego h[i].contraparte = u}
    }

    \bigskip

    \predLargo{todosExclusivos}{usuarios: \dict{\Tipo{Id}}{\Tipo{Historial}}, tuplas: \seq{\tupla{\Tipo{Id}, \Tipo{Id}}}}{
        \paraTodo{i}{\Z}{0 \leq i < |tuplas| \implicaLuego \\
            \tab (
                tuplas[i]_0 \in claves(usuarios) \land 
                tuplas[i]_1 \in claves(usuarios) \yLuego \\
                \tab\tab sonExclusivos(usuarios[tuplas[i]_0], usuarios[tuplas[i]_1], tuplas[i]_0, tuplas[i]_1)
            )
        }
    }

\bigskip

\predLargo{todosLosExclusivosEstan}{usuarios: \dict{\Tipo{Id}}{\Tipo{Historial}}, tuplas: \seq{\tupla{\Tipo{Id}, \Tipo{Id}}}}{
    \paraTodo{u_1, u_2}{\Tipo{Id}}{(
        u_1 \in claves(usuarios) \land u_2 \in claves(usuarios) \land u_1 \neq u_2 \land \\
        \tab sonExclusivos(usuarios[u_1], usuarios[u_2], u_1, u_2)
    ) \implicaLuego (
        \existe{i}{\Z}{0 \leq i < |tuplas| \yLuego (tuplas[i] = \tupla{u_1, u_2} \lor tuplas[i] = \tupla{u_2, u_1})}
    )}
}

\bigskip

\predLargo{sinEspejosEnTuplas}{tuplas: \seq{\tupla{\Tipo{Id}, \Tipo{Id}}}}{
    \neg \existe{i,j}{\Z}{0 \leq i < |tuplas| \land 0 \leq j < |tuplas| \land i \neq j \yLuego \\
        \tab (tuplas[i]_0 = tuplas[j]_1 \land tuplas[i]_1 = tuplas[j]_0)
    }
}

\bigskip

\subsection{Procedimientos}

\bigskip

    \begin{proc}{crearSistema}{}
        {\Tipo{Airmiles}}
       
        \requiere{True}
        \aseguraLargo{
            res.usuarios = \{ \ \} \ \land \\ 
            res.promos = \{ "NoPromo":\struct{1,0,False}\}
        }    
        
    \end{proc} \\

    \begin{proc}{registrarUsuario}{\Inout a: \Tipo{AirMiles}}{}
    \requiere{a = A_0}
    \aseguraLargo{
        \existe{k}{Id}{k > 0 \land k \notin claves(A_0.usuarios) \yLuego \\
        a.usuarios = setKey(A_0.usuarios,k,\langle \rangle)} \land \\
        a.promos = A_0.promos
    }
\end{proc}

\bigskip

\begin{proc}{registrarUsuarioConReferidor}{\Inout a: \Tipo{AirMiles}, \In u: \Tipo{Id}}{}
    \requiere{a = a_0 \land u \in claves(a_0.usuarios)}
    \aseguraLargo{
        a.promos = a_0.promos \land \\
        \tab (\exists i: Id)\ ( \\
        \tab\tab i > 0 \land i \notin claves(a_0.usuarios) \yLuego \\
        \tab\tab a.usuarios = setKey( \\
        \tab\tab\tab setKey( \\
        \tab\tab\tab\tab a_0.usuarios, \\
        \tab\tab\tab\tab i, \\
        \tab\tab\tab\tab \langle \langle \Tipo{REFERIDO}, 0, 500, a_0.promos[NoPromo].factor, \textnormal{''NoPromo''}, u \rangle \rangle), \\
        \tab\tab\tab u, \\
        \tab\tab\tab a_0.usuarios[u] \concat \langle \langle \Tipo{REFERIDO}, 0, 500, a_0.promos[NoPromo].factor, \textnormal{''NoPromo''}, i \rangle \rangle \\
        \tab\tab ) \\
        \tab )
    }
\end{proc}

\bigskip

% CORRECCIÓN: "El enunciado pide guardar distancia volada y factor de promoción aplicado."
\begin{proc}{acumularMillas}{\Inout a: \Tipo{AirMiles}, \In Id: \Tipo{Id}, \In distancia: \R, \In codigo: \Tipo{Codigo}}{}
    \requiereLargo{
        a = A_0 \land distancia > 0 \land Id \in claves{a.usuarios} \land codigo \in claves(a.promos) \yLuego \\
        esPromoValida(\\
            \tab codigo\\
        \tab millasGanadasPorVuelo(distancia, A_0.usuarios[Id], A_0.promos[codigo].factor) \\
        )
    }
    \aseguraLargo{
        a.usuarios = setKey( \\
        \tab A_0.usuarios, \\
        \tab Id, \\
        \tab A_0.usuarios[Id] ++ \\
        \tab\tab \langle \langle \\
        \tab\tab\tab \Tipo{VUELO}, \\
        \tab\tab\tab distancia, \\
        \tab\tab\tab millasGanadasPorVuelo( \\
        \tab\tab\tab\tab distancia, \\
        \tab\tab\tab\tab A_0.usuarios[Id], \\
        \tab\tab\tab\tab A_0.promos[codigo].factor \\
        \tab\tab\tab ), \\
        \tab\tab\tab A_0.promos[codigo].factor, \\
        \tab\tab\tab codigo, \\
        \tab\tab\tab sinContraparte() \\
        \tab\tab \rangle \rangle \\
        ) \land ( \\
        \tab (\lnot A_0.promos[codigo].tieneTope \implica a.promos = A_0.promos) \land \\
        \tab (A_0.promos[codigo].tieneTope \land \\
        \tab\tab millasGanadasPorVuelo( \\
        \tab\tab\tab distancia, \\
        \tab\tab\tab A_0.usuarios[Id], \\
        \tab\tab\tab A_0.promos[codigo].factor \\
        \tab\tab ) < A_0.promos[codigo].tope \implica \\
        \tab\tab\tab a.promos = setKey( \\
        \tab\tab\tab\tab A_0.promos, \\
        \tab\tab\tab\tab codigo, \\
        \tab\tab\tab\tab \langle \\
        \tab\tab\tab\tab\tab A_0.promos[codigo].factor, \\
        \tab\tab\tab\tab\tab A_0.promos[codigo].tope - \\
        \tab\tab\tab\tab\tab\tab millasGanadasPorVuelo( \\
        \tab\tab\tab\tab\tab\tab\tab distancia, \\
        \tab\tab\tab\tab\tab\tab\tab A_0.usuarios[Id], \\
        \tab\tab\tab\tab\tab\tab\tab A_0.promos[codigo].factor \\
        \tab\tab\tab\tab\tab\tab ), \\
        \tab\tab\tab\tab\tab A_0.promos[codigo].tieneTope \\
        \tab\tab\tab\tab \rangle \\
        \tab\tab\tab ) \\
        \tab ) \land \\
        \tab (A_0.promos[codigo].tieneTope \land \\
        \tab\tab millasGanadasPorVuelo( \\
        \tab\tab\tab distancia, \\
        \tab\tab\tab A_0.usuarios[Id], \\
        \tab\tab\tab A_0.promos[codigo].factor \\
        \tab\tab ) = A_0.promos[codigo].tope) \implica \\
        \tab\tab\tab\tab a.promos = delKey(A_0.promos, codigo) \\
        \tab ) \\
        
        
    }
\end{proc}

\newpage

% CORRECCIÓN: "Por lo general, por convención se pone primero el inout."
% (asegura): se reescribió para el nuevo historial; ya usa A0.usuarios, como pidió el corrector.
\begin{proc}{canjearMillas}{
    
    \Inout a: \Tipo{AirMiles},
    \In cantidadMillas: \Tipo{Millas}, 
    \In u: \Tipo{Id}, 
    \In premio: \Tipo{Codigo} 
    }{}
    \requiere{
        a = A_0\ \land\ cantidadMillas > 0\ \land\ u \in a.usuarios\ \yLuego \\
        \tab\tab millasDisponibles(a.usuarios[u]) \geq cantidadMillas
    }
    \aseguraLargo{ 
        a.usuarios = \mathsf{setKey}(A_0.usuarios, u, \\
        A_0.usuarios[u] \concat \langle \langle CANJE, 0, cantidadMillas,1, premio, sinContraparte() \rangle \rangle) \land \\
        a.promos = A_0.promos
    }
\end{proc}

\bigskip

% ---------------------------------------------------------------------
% transferirMillas (Héctor)
% Correcciones aplicadas:
%  - requiere: el yLuego ahora protege a todas las indexaciones (después de validar u1, u2 y u1 != u2);
%              se sacó el primer yLuego, que no hacía falta.
%  - requiere: factorCategoriaUsuario se aplica al historial de u_2, no al Id.
%  - asegura:  se usa siempre A_0 (antes aparecían a_0, a1 y u, que no estaban definidos).
%  - asegura:  ya no se agrega una secuencia <<u2>> a emision: la secuencia emision se eliminó.
%  - asegura:  "El receptor no guarda de quién recibió" -> ahora cada transacción guarda su contraparte:
%              el emisor guarda a u_2 (a quién envió) y el receptor guarda a u_1 (de quién recibió).
% ---------------------------------------------------------------------
\begin{proc}{transferirMillas}{\Inout a: \Tipo{AirMiles},\In u_1: \Tipo{Id},\In u_2: \Tipo{Id},\In m: \Tipo{Millas}}
    {}
    \requiereLargo{
        (a = A_0 \land u_1 \in claves{A_0.usuarios} \land u_2 \in claves{A_0.usuarios} \land u_1 \neq u_2) \yLuego \\
        (m > 0 \land \\
        \tab millasDisponibles(A_0.usuarios[u_1]) \geq m \land \\
        \tab factorCategoriaUsuario(A_0.usuarios[u_2]) \neq 0)
    }
    \aseguraLargo{
        a.promos = A_0.promos \land \\
        a.usuarios = setKey( \\
        \tab setKey( \\
        \tab\tab A_0.usuarios, \\
        \tab\tab u_1, \\
        \tab\tab A_0.usuarios[u_1] \concat \langle \langle \Tipo{EMISOR}, 0, m, a_0.promos[NoPromo].factor , \comillas{NoPromo}, u_2 \rangle \rangle \\
        \tab ), \\
        \tab u_2, \\
        \tab A_0.usuarios[u_2] \concat \langle \langle \Tipo{RECEPTOR}, 0, m, a_0.promos[NoPromo].factor , \comillas{NoPromo}, u_1 \rangle \rangle )
    }
\end{proc}

\bigskip

\begin{proc}{consultarSaldo}{\In a: \Tipo{AirMiles}, \In u: \Tipo{Id}}{\Tipo{Millas}}
    \requiere{u \in claves(a.usuarios)}
    \asegura{
        res = millasDisponibles(a.usuarios[u])
    }
\end{proc}

\bigskip

\begin{proc}{consultarCategoria}
{
\In a: \Tipo{AirMiles},
\In u: \Tipo{Id}
}{\Tipo{Categoria}}
    \requiere{u \in claves(a.usuarios)}
    \aseguraLargo{
        (factorCategoriaUsuario(a.usuarios[u]) = 1 \implica res = \text{BRONCE}) \land \\
        (factorCategoriaUsuario(a.usuarios[u]) = 1.5 \implica res = \text{PLATA}) \land \\
        (factorCategoriaUsuario(a.usuarios[u]) = 2 \implica res = \text{ORO}) \land \\
        (factorCategoriaUsuario(a.usuarios[u]) = 0 \implica res = \text{PLATINO})
    }
\end{proc}

\newpage

    \begin{proc}{rankingMillas}{
        \In a:\Tipo{AirMiles},
        \In n:\Z 
        }
        {\seq{\Tipo{Id}}
        }
        \requiere{ \ n \geq 1}
        \aseguraLargo{
            (\,|res| = n \land n < |a.usuarios|\,) 
            \oplus 
            (\,|res| = |a.usuarios| \land n \geq |a.usuarios|\,) \\
            ~ \land ~ esRankingOrdenado(res, a.usuarios) \\ ~ \land ~ tieneAlosTopN(res,a.usuarios) \\
            ~ \land ~ sonTodosUsuarios(ranking, a.usuarios) \\ \land ~ sinRepetidos(ranking)
        }
    \end{proc} \\  

\bigskip

% ---------------------------------------------------------------------
% resumenTransacciones (Héctor)
% Corregido: devuelve 4 valores para no mezclar Millas enviadas y recibidas:
% <por VUELO, canjeadas, transferidas enviadas, transferidas recibidas>
% ---------------------------------------------------------------------
\begin{proc}{resumenTransacciones}{\In u: \Tipo{Id}, \In a: \Tipo{AirMiles}}{\langle \R, \R, \R, \R \rangle}
    \requiere{u \in claves(a.usuarios)}
    \aseguraLargo{
        res = \langle totalMillasPorVuelo(a.usuarios[u]), \\
        \tab totalMillasCanjeadas(a.usuarios[u]), \\
        \tab totalMillasEmitidas(a.usuarios[u]), \\
        \tab totalMillasRecibidas(a.usuarios[u]) \rangle
    }
\end{proc}

\bigskip

\bigskip
\begin{proc}{iniciarPromocion}
    {
    \Inout a: \Tipo{AirMiles},
    \In codigo: \Tipo{Codigo},
    \In factor: \Tipo{Factor}
    }{}
        \requiereLargo{
            a = a_0 \land \\
            codigo \notin claves(a.promos) \land \\
            factor \geq 1
        }
        \aseguraLargo{
            a.usuarios = a_0.usuarios \land \\
            a.promos = setKey(
            a_0.promos,
            codigo, 
            \tupla{Factor, 0, False}
            )
        }
    \end{proc}\\

\bigskip

    \begin{proc}{iniciarPromocionConTope}
    {
    \Inout a: \Tipo{AirMiles},
    \In codigo: \Tipo{Codigo},
    \In factor: \Tipo{Factor},
    \In Tope: \Tipo{Tope}
    }{}
        \requiereLargo{
            a = a_0 \land \\
            factor \geq 1 \land \\
            Tope > 0 \land \\
            codigo \notin claves(a.promos)
        }
        \aseguraLargo{
            a.usuarios = a_0.usuarios \land \\
            a.promos = setKey(
            a_0.promos,
            codigo, 
            \tupla{Factor, Tope, True}
            )
        }
    \end{proc}

\newpage

    \begin{proc}{finalizarPromocion}{
        \Inout a: \Tipo{AirMiles},
        \In promoActiva: \Tipo{Codigo} 
        }
        {}

        \requiere{ a = A_0 ~ \land ~ promoActiva \in a.promos ~ \land ~ \neg (\,promoActiva = "NoPromo") }
        \aseguraLargo{
            a.usuarios = A_0.usuarios ~ \land ~
            a.promos = delKey(A_0.promos,promoActiva)
        }
    \end{proc} \\

\bigskip

% ---------------------------------------------------------------------
% elQueTieneMasMillas (Héctor)
% CORRECCIÓN: "Innecesario el existe res. Ya tenemos res como palabra reservada
%              para lo que devuelve el proc."
% ---------------------------------------------------------------------
\begin{proc}{elQueTieneMasMillas}{\In a: \Tipo{AirMiles}}{Id}
    \requiere{|claves(a.usuarios)| \geq 1}
    \aseguraLargo { (res \in claves(a.usuarios) \yLuego \\
    (\forall k: \Tipo{Id})\ (k \in claves(a.usuarios) \implicaLuego \\
     millasAcumuladas(a.usuarios[k]) \leq 
     millasAcumuladas(a.usuarios[res]))) }
\end{proc}

\bigskip

\begin{proc}{elMasCanjeador}
{
\In a: \Tipo{AirMiles}
}{\Tipo{Id}}
    \requiere{|claves(a.usuarios)| \geq 1}
    \aseguraLargo{
        \existe{i}{\Tipo{Id}}{
            i \in claves(a.usuarios) \land_l \\
            \paraTodo{u}{\Tipo{Id}}{
                u \in claves(a.usuarios) \implicaLuego \\
                sumaDeMillas(a.usuarios[i], \langle CANJE, EMISOR \rangle) \geq \\ sumaDeMillas(a.usuarios[u], \langle CANJE, EMISOR \rangle)
            } \\ \land res = i
        }
    }
\end{proc}\\

% \begin{proc}{paresTransferidoresExclusivos}
% {
% \In a: \Tipo{AirMiles}
% }{\seq{\tupla{\Tipo{Id}, \Tipo{Id}}}}
%     \requiere{True}
%     \aseguraLargo{
%         \paraTodo{u1,u2}{\Tipo{Id}}{ \\
%             (u1 \in claves(a.usuarios) \land u2 \in claves(a.usuarios) \land u1 \neq u2) \yLuego \\
%             \tab sonExclusivos(a.usuarios[u1].emisiones, a.usuarios[u2].emisiones, u1, u2 \\)
%             \implicaLuego \\
%             \tab \existe{i}{\Z}{
%                 0 \leq i < |res| \yLuego res[i] = \tupla{u1,u2} \land \\
%                 \tab\tab \paraTodo{k}{\Z}{
%                     0 \leq k < |res| \yLuego res[k] = \tupla{u1,u2} \implicaLuego k = i
%                 }
%             } \land \\
%             \tab\lnot \existe{j}{\Z}{0 \leq j < |res| \yLuego res[j] = \tupla{u2,u1}} \\
%             \tab\land \neg(\exists <us1, us2> : tupla\langle\Z,\Z\rangle) \\
%             \tab(<us1, us2> \in res \land ((us1,us2 \notin a.usuarios) \lor \\
%             \tab\tab(us1, us2 \in a.usuarios \land_l\\ 
%             \tab\tab \neg sonExclusivos(a.usuarios[us1], 
%             \\ \tab\tab\tab\tab\tab\tab a.usuarios[us2], 
%             \\ \tab\tab\tab\tab\tab\tab us1, us2))))
%         }
%     }
% \end{proc}\\


% \begin{proc}{paresTransferidoresExclusivos}
% {
% \In a: \Tipo{AirMiles}
% }{\seq{\tupla{\Tipo{Id}, \Tipo{Id}}}}
%     \requiere{True}
%     \aseguraLargo{
%         \paraTodo{u1,u2}{\Tipo{Id}}{ \\
%             (u1 \in claves(a.usuarios) \land u2 \in claves(a.usuarios) \land u1 \neq u2) \yLuego \\
%             \tab sonExclusivos(a.usuarios[u1].emisiones, a.usuarios[u2].emisiones, u1, u2 \\)
%             \implicaLuego esUnicaTuplaDeElementos(u1, u2, res) ~ \land ~ sonParesExclusivos(res, usuarios)
%         }
%     }
% \end{proc}\\

% \begin{proc}{paresTransferidoresExclusivos}
% {
% \In a: \Tipo{AirMiles}
% }{\seq{\tupla{\Tipo{Id}, \Tipo{Id}}}}
%     \requiere{True}
%     \aseguraLargo{
%         \paraTodo{u1,u2}{\Tipo{Id}}{ \\
%             (u1 \in claves(a.usuarios) \land u2 \in claves(a.usuarios) \land u1 \neq u2) \yLuego \\
%             \tab sonExclusivos(a.usuarios[u1].emisiones, a.usuarios[u2].emisiones, u1, u2 \\)
%             \implicaLuego esUnicaTuplaDeElementos(u1, u2, res) ~ \land ~ sonParesExclusivos(res, usuarios)
%         }
%     }
% \end{proc}\\

\begin{proc}{paresTransferidoresExclusivos}
{\In a: \Tipo{AirMiles}}{\seq{\tupla{\Tipo{Id}, \Tipo{Id}}}}
    \requiere{True}
    \aseguraLargo{
        todosExclusivos(a.usuarios, res) \land \\
        todosLosExclusivosEstan(a.usuarios, res) \land \\
        sinParesRepetidos(res) \land \\
        sinEspejosEnTuplas(res)
    }
\end{proc}


% \predLargo{esUnicaTuplaDeElementos}{
%     el1: \Tipo{id},
%     el2:\Tipo{id},
%     tuplas: \seq{\tupla{\Tipo{id},\Tipo{id}}}
%     }
%     {
%         \tab \existe{i}{\Z}{
%             0 \leq i < |tuplas| \yLuego tuplas[i] = \tupla{el1,el2} \land \\
%             \tab\tab \paraTodo{k}{\Z}{
%                 0 \leq k < |tuplas| \yLuego tuplas[k] = \tupla{el1,el2} \implicaLuego k = i
%             }
%         } \land \\
%         \tab\lnot \existe{j}{\Z}{0 \leq j < |tuplas| \yLuego res[j] = \tupla{el2,el1}} \\
%     }

% \bigskip

% \predLargo{sonParesExclusivos}{pares:\seq{\tupla{\Tipo{id},\Tipo{id}}}, usuarios: \dict{\Tipo{Id}}{\Tipo{Historial}}}{
%     \neg(\exists <u1, u2> : tupla\langle\Z,\Z\rangle) \\
%     \tab(<u1, u2> \in pares \land \\ 
%     \tab\tab(u1,u2 \notin usuarios \lor \\
%     \tab\tab\tab(
%         u1, u2 \in usuarios ~ \yLuego ~
%         \neg sonExclusivos(usuarios[u1], usuarios[u2], u1, u2)
%         )\\
%     \tab\tab)\\
%     \tab)
% }


% % Propuesta de claudio
% \begin{proc}{paresTransferidoresExclusivos}{\In a: \Tipo{AirMiles}}{\seq{\langle \Tipo{id}, \Tipo{id} \rangle}}
%     \requiereCorto{True}
%     \asegura{
%         sinRepetidos(res) \land \\
%         todosSonExclusivos(res, a.usuarios) \land \\
%         estanTodosLosExclusivos(res, a.usuarios) \land \\
%         sinEspejos(res)
%     }
% \end{proc}

% % ----- Predicados de paresTransferidoresExclusivos -----
% % Para las tuplas se usa la notación de la materia: s[i]_0 y s[i]_1 son la primera y la segunda componente.

% % u_1 y u_2 son dos usuarios distintos del sistema (protege las indexaciones en usuarios).
% \pred{sonUsuariosDistintos}{u_1: \Tipo{id}, u_2: \Tipo{id}, usuarios: \dict{\Tipo{id}}{\Tipo{historial}}}{
%     u_1 \in \claves{usuarios} \land u_2 \in \claves{usuarios} \land u_1 \neq u_2
% }

% % 1) No hay tuplas repetidas en s.
% \pred{sinRepetidos}{s: \seq{\langle \Tipo{id}, \Tipo{id} \rangle}}{
%     (\forall i, j: \Z)\ ((0 \leq i < |s| \land 0 \leq j < |s| \land i \neq j) \implicaLuego s[i] \neq s[j])
% }

% % 2) Toda tupla de s es un par de usuarios exclusivos.
% \pred{todosSonExclusivos}{s: \seq{\langle \Tipo{id}, \Tipo{id} \rangle}, usuarios: \dict{\Tipo{id}}{\Tipo{historial}}}{
%     \paraTodo{i}{\Z}{0 \leq i < |s| \implicaLuego \\
%     \tab sonUsuariosDistintos(s[i]_0, s[i]_1, usuarios) \yLuego \\
%     \tab sonExclusivos(usuarios[s[i]_0], usuarios[s[i]_1], s[i]_0, s[i]_1)}
% }

% % 3) Todo par de usuarios exclusivos del observador usuarios aparece en s (en alguno de los dos órdenes).
% \pred{estanTodosLosExclusivos}{s: \seq{\langle \Tipo{id}, \Tipo{id} \rangle}, usuarios: \dict{\Tipo{id}}{\Tipo{historial}}}{
%     (\forall u_1, u_2: \Tipo{id})\ ( \\
%     \tab (sonUsuariosDistintos(u_1, u_2, usuarios) \yLuego \\
%     \tab\ sonExclusivos(usuarios[u_1], usuarios[u_2], u_1, u_2)) \\
%     \tab \implicaLuego (\langle u_1, u_2 \rangle \in s \lor \langle u_2, u_1 \rangle \in s) \\
%     )
% }

% % 4) No hay tuplas espejo en s: si <u_1, u_2> está en s, entonces <u_2, u_1> no está.
% \pred{sinEspejos}{s: \seq{\langle \Tipo{id}, \Tipo{id} \rangle}}{
%     (\forall i, j: \Z)\ ((0 \leq i < |s| \land 0 \leq j < |s|) \implicaLuego s[i] \neq \langle s[j]_1, s[j]_0 \rangle)
% }

\bigskip


\end{tad}

\end{document}
